\subsection{Vector spaces over an ordered field}

Let K be an ordered field, and let E be a vector space over K. The relation « there exists \(\lambda > 0\) in K such that \(y = \lambda x \gg\) between two elements x and y in the set \(E - \{0\}\) is an equivalence relation. The equivalence classes under this relation are called open half-lines with origin 0; the union of an open half-line and \(\{0\}\) is called a closed half-line (or sometimes simply a half-line) with origin 0. Every vector \(a \neq 0\) contained in an open (resp. closed) half-line \(\Delta\) is called a direction vector of \(\Delta\) , and \(\Delta\) is the set of vectors \(\lambda a\) for all scalars \(\lambda > 0\) (resp. \(\lambda \geqslant 0\) ). Every line D through 0 contains exactly two open (resp. closed) half-lines with origin 0; if \(\Delta\) is one of these, then \(-\Delta\) is the other (called the opposite of \(\Delta\) ).

Now if F is an affine space over K, and E the space of translations of F, then any subset of F of the form \(\Delta = a + \Delta_{0}\) , where \(\Delta_{0}\) is an open (resp. closed) half-line of E, is called an open (resp. closed) half-line with origin \(a \in F\) . The half-line \(\Delta_{0}\) is completely determined by \(\Delta\) (for it is the half-line with direction vector b - a, for any \(b \neq a\) in \(\Delta\) ), and is called the direction of \(\Delta\) ; a direction vector of \(\Delta_{0}\) is also called a direction vector of \(\Delta\) .

Suppose now that E has finite dimension n over K; then it is known (III, p.~518, Cor. 1) that the n-th exterior power \(\bigwedge^{n}E\) is a vector space of dimension 1 over K, and so the union of two opposite closed half-lines of origin 0. These half-lines are called orientations of E; the space E together with a given one of these half-lines \(\Delta\) is said to be oriented; an \(n\) -vector \(z\) is then called positive (resp. negative) under this orientation if it belongs to \(\Delta\) (resp. to \(-\Delta\) ); it is negative (resp. positive) under the opposite orientation.

An orientation of an affine space F over K is by definition an orientation of the space of translations of F; the space F, together with such an orientation, is called an oriented affine space.

Let \(\mathbf{E}\) be an oriented vector space over \(\mathbf{K}\) , of dimension \(n\) ; an ordered basis \((a_i)_{1 \leqslant i \leqslant n}\) of \(\mathbf{E}\) is called positive or direct (resp. negative or inverse) if the \(n\) -vector \(a_1 \wedge a_2 \wedge \ldots \wedge a_n\) is positive (resp. negative). If \(u\) is an automorphism of the vector space \(\mathbf{E}\) then \((\bigwedge^{n} u)(z) = \det(u) \cdot z\) for all \(z \in \bigwedge^{n} \mathbf{E}\) , so a necessary and sufficient

condition for \(\bigwedge^{n} u\) to leave invariant the orientation of E (or, as we also say, to preserve the orientation) is that \(\det(u) > 0\) ; the automorphisms having this property are precisely the automorphisms of E as an oriented vector space; they form a normal subgroup \(\mathbf{GL}^{+}(\mathbf{E})\) of the linear group \(\mathbf{GL}(\mathbf{E})\) , which has index 2 whenever \(E \neq 0\) .

When \(\mathbf{E} = 0\) then \(\mathbf{GL}(\mathbf{E}) = \operatorname{End}(\mathbf{E})\) contains only the identity map \(1_{\mathbf{E}}\) and by definition \(\det(1_{\mathbf{E}}) = 1\) . Note that \(\bigwedge^{n}\mathbf{E} = \bigwedge^{0}\mathbf{E} = \mathbf{K}\) by definition in this case; the half-line of \(\mathbf{K}\) formed by the positive scalars is called the canonical orientation of the zero space.

Let \(\mathbf{M}\) and \(\mathbf{N}\) be two complementary subspaces of dimensions \(p\) and \(n - p\) respectively in the vector space \(\mathbf{E}\) of dimension \(n\) ; if \(z'\) (resp. \(z''\) ) is a nonzero vector in \(\bigwedge^{p} \mathbf{M}\) (resp. \(\bigwedge^{n - p} \mathbf{N}\) ), then \(z' \wedge z''\) is a nonzero vector in

\(\bigwedge^{n} \mathbf{E}\). Given an orientation on M and an orientation on N, the vectors \(z' \wedge z''\) for positive \(z'\) and \(z''\) form an orientation of E, called the product orientation of the orientation of M by the orientation of N (which depends on the order of the factors when \(p(n - p)\) is odd). Conversely, given orientations on E and on M, there exists a unique orientation on N such that the given orientation on E is the product of the given orientation on M and this orientation on N (in that order); this orientation is said to be complementary to the orientation of M with respect to that of E. If \(N'\) is a second complementary subspace to M, the canonical projection \(N \to N'\) parallel to M takes the complementary orientation of N onto that of \(N'\) . The image of the complementary orientation of N under the canonical map \(N \to E/M\) is thus independent of the choice of complement N; it is called the quotient orientation on E/M of the orientation of E by that of M.

